\chapter{Luteinizing Hormone (LH)}

\section*{What Is This Test?}
Luteinizing hormone (LH, also called lutropin) is a glycoprotein hormone produced by gonadotrope cells of the anterior pituitary gland. Like FSH, LH consists of a common alpha subunit (shared with FSH, TSH, and hCG) and a unique beta subunit that confers receptor specificity. LH is secreted in a pulsatile manner under the control of gonadotropin-releasing hormone (GnRH) from the hypothalamus, with pulse frequency and amplitude determining the pattern of LH secretion. In females, LH acts on ovarian theca cells to stimulate androgen (androstenedione and testosterone) production, which is then aromatized to estradiol by granulosa cells under FSH stimulation (the two-cell, two-gonadotropin model). The mid-cycle LH surge (a 10--20 fold increase sustained for 36--48 hours) triggers ovulation approximately 36 hours after the surge onset, luteinization of the ruptured follicle, and progesterone production by the corpus luteum. In males, LH acts on testicular Leydig cells to stimulate testosterone production. LH secretion is under negative feedback from estradiol and testosterone. In women, estradiol switches from negative to positive feedback during the late follicular phase (estradiol >200 pg/mL for >36 hours), generating the pre-ovulatory LH surge. The LH test is used to evaluate ovulation, classify amenorrhea, diagnose PCOS, assess male hypogonadism, evaluate pubertal development, and monitor ovulation induction.

\section*{Alternative Names}
Lutropin; Serum LH; Plasma LH; Interstitial Cell-Stimulating Hormone (ICSH, historical); LH:FSH Ratio; LH Surge; Ovulation Predictor

\section*{Principle}
LH is measured using a two-site sandwich immunometric (chemiluminescent) immunoassay on automated analyzers. Two monoclonal antibodies specific for the beta subunit of LH (with minimal cross-reactivity to FSH, TSH, or hCG, though there is approximately 1--5\% cross-reactivity with hCG due to structural similarity) are used: a capture antibody on a solid phase and a detection antibody labeled with a chemiluminescent tag (ruthenium on Roche ECLIA, acridinium ester on Abbott/Siemens). After incubation and washing, the signal is proportional to LH concentration. Third-generation assays have functional sensitivity of <0.1 mIU/mL and are calibrated against the WHO 2nd International Standard (IS 80/552) \cite{tietz2023textbook}. Inter-assay CV is <5\%. LH is more pulsatile than FSH, with pulses occurring every 60--120 minutes in the follicular phase and every 60--90 minutes in the luteal phase. A single LH measurement captures a random point in this pulsatile secretion, which is important for interpretation. The same assay and laboratory should be used for serial monitoring.

\section*{History}
LH was identified as a distinct gonadotropin in the 1930s. The term luteinizing hormone was coined because of its effect on corpus luteum formation in the rat ovary. The first bioassay for LH (the ovarian ascorbic acid depletion assay) was developed by Parlow in 1958. The first radioimmunoassay was developed by Odell and colleagues in 1967 \cite{odell1967radioimmunoassay}. The mid-cycle LH surge as the trigger for ovulation was first characterized in the 1970s. The development of urinary LH kits for home ovulation prediction in the 1980s revolutionized fertility awareness. The role of elevated LH in PCOS pathophysiology was elucidated in the 1980s--1990s, and the LH:FSH ratio >2:1 was incorporated into PCOS diagnostic criteria, though later de-emphasized in the Rotterdam criteria (2003) in favor of hyperandrogenism, oligo-ovulation, and polycystic ovarian morphology \cite{rotterdam2004revised}. The pulsatile nature of LH secretion was characterized by Knobil and colleagues in the 1970s using rhesus monkeys, establishing the critical role of GnRH pulse frequency in regulating gonadotropin secretion \cite{knobil1980neuroendocrine}.

\section*{Why This Test Is Ordered}
\begin{itemize}
  \item Evaluation of ovulation and menstrual cycle function --- mid-cycle LH surge confirms ovulatory cycles; absent LH surge indicates anovulation
  \item Diagnosis of PCOS --- LH:FSH ratio >2:1 is a supportive finding, present in approximately 60\% of PCOS patients \cite{balen2016management}
  \item Classification of hypogonadism in men: elevated LH with low testosterone indicates primary (testicular) hypogonadism; low or inappropriately normal LH with low testosterone indicates secondary (pituitary/hypothalamic) hypogonadism \cite{bhasin2018testosterone}
  \item Diagnosis of menopause (elevated LH >30 mIU/mL with elevated FSH) and premature ovarian insufficiency
  \item Monitoring ovarian stimulation during IVF cycles --- LH levels guide gonadotropin dosing and timing of hCG trigger
  \item Assessment of pubertal development --- elevated LH in children <8 years (girls) or <9 years (boys) indicates central precocious puberty (LH response >5--8 mIU/mL to GnRH stimulation is diagnostic)
  \item Investigation of suspected pituitary dysfunction --- low LH with low testosterone or estradiol suggests gonadotropin deficiency
  \item Evaluation of delayed puberty --- distinguish constitutional delay (low LH for chronologic age but appropriate for bone age) from permanent hypogonadotropic hypogonadism
  \item Home ovulation prediction using urinary LH surge detection kits
\end{itemize}

\section*{Is This a Primary or Secondary Test}
LH is a primary test for evaluating ovulation function, classifying male hypogonadism, and diagnosing menopause/POI. It is part of the standard infertility and hypogonadism evaluations, always interpreted alongside FSH, estradiol (women), and testosterone (men). In PCOS, the LH:FSH ratio is a supportive (secondary) finding.

\section*{How This Test Is Performed}
A healthcare professional draws a blood sample from a vein into a serum separator tube (gold-top) or plain red-top tube (~3--5 mL). In menstruating women evaluated for PCOS or infertility, LH should be drawn in the early follicular phase (day 2--4), as LH is lowest at this time and the LH:FSH ratio is most informative \cite{mihm2011normal}. For ovulation detection, serial samples across the mid-cycle (days 10--16) are needed. In men, LH should be drawn in the morning (7--10 AM) alongside testosterone. LH is stable at room temperature for 24 hours and refrigerated for 7 days.

\section*{What Preparation Is Needed}
No fasting is required. Cycle day must be documented. Oral contraceptives and estrogen therapy suppress LH and should be held 4--6 weeks before testing if evaluating pituitary function. Exogenous testosterone and anabolic steroids suppress LH. GnRH agonists/antagonists profoundly alter LH. High-dose biotin (>5 mg/day) should be held 48 hours before testing. Opiates suppress LH. Clomiphene and tamoxifen elevate LH. Stress, acute illness, and strenuous exercise can alter LH secretion.

\section*{Contraindications and Precautions}
LH is highly pulsatile (pulses every 60--120 minutes), and a single random value may not represent the integrated LH secretion. For clinical classification, FSH (less pulsatile, longer half-life) is often more reliable. LH values MUST be interpreted in the context of menstrual cycle phase, patient age, sex, and concurrent medications. In perimenopause, LH fluctuates widely. In PCOS, LH is typically elevated with increased pulse amplitude and frequency. hCG cross-reacts in some LH assays (~1--5\%), causing minor elevation in early pregnancy that should not be confused with an LH surge. Chronic kidney disease elevates LH through decreased clearance.

\section*{Risks}
Minimal risk. Slight pain, bruising, or hematoma at venipuncture site. Rarely: vasovagal syncope.

\section*{What Happens After the Test}
Results available within 1--4 hours. Interpretation with FSH and sex steroids is essential. In women: elevated LH with normal FSH and elevated LH:FSH ratio supports PCOS. Very elevated LH (>30 mIU/mL) with elevated FSH suggests menopause or POI. In men: elevated LH with low testosterone indicates primary hypogonadism (Leydig cell failure). Low LH with low testosterone indicates secondary hypogonadism requiring pituitary MRI.

\section*{Understanding Results}

\subsection*{Below Normal}
\begin{itemize}
  \item Secondary (hypogonadotropic) hypogonadism: low LH with low testosterone (male) or low estradiol (female). Causes include functional hypothalamic amenorrhea, Kallmann syndrome, pituitary adenomas/tumors, hyperprolactinemia, chronic opioid use, and constitutional delay of puberty.
  \item Anovulation: absence of mid-cycle LH surge. Common in PCOS, hypothalamic amenorrhea, and hyperprolactinemia. Confirmed by low luteal phase progesterone.
  \item Oral contraceptive use: LH suppressed to <5 mIU/mL.
  \item GnRH agonist/antagonist therapy: LH suppressed in IVF cycles and prostate cancer treatment.
  \item Isolated LH deficiency (fertile eunuch syndrome): rare, low LH with normal FSH causing hypogonadism with preserved spermatogenesis.
\end{itemize}

\subsection*{Above Normal}
\begin{itemize}
  \item Menopause and premature ovarian insufficiency: elevated LH (>30 mIU/mL, often 50--200 mIU/mL) with elevated FSH. Loss of ovarian estradiol and inhibin negative feedback causes unopposed GnRH-driven gonadotropin secretion. LH rises later and to a lesser degree than FSH.
  \item PCOS: elevated LH with normal or low-normal FSH, LH:FSH ratio >2:1. Elevated LH pulse frequency and amplitude due to increased GnRH pulse generator activity from hyperandrogenemia and relative progesterone deficiency. Present in ~60\% of PCOS.
  \item Primary testicular failure: elevated LH (often >15--20 mIU/mL) with low testosterone. Klinefelter syndrome (47,XXY) is the most common cause. LH rises as testosterone negative feedback is lost.
  \item Mid-cycle LH surge: LH peaks at 25--100 mIU/mL (10--20 fold above baseline), sustained for 36--48 hours. Used to time intercourse or IUI (ovulation occurs ~36 hours after surge onset). Urinary LH kits detect the surge.
  \item Gonadal dysgenesis in children: elevated gonadotropins in Turner syndrome (45,X) or Swyer syndrome (46,XY).
  \item GnRH agonist flare response: initial agonistic stimulation causes transient LH elevation before receptor desensitization and suppression. Seen in IVF and prostate cancer treatment.
\end{itemize}

\section*{Normal Results}
\begin{itemize}
  \item Adult female, follicular phase (day 2--4): 2--10 mIU/mL
  \item Adult female, mid-cycle surge peak: 25--100 mIU/mL
  \item Adult female, luteal phase: 1--8 mIU/mL
  \item Adult female, postmenopausal: 15--60 mIU/mL
  \item Adult male (18--60 years): 1.5--9 mIU/mL
  \item Adult male (>60 years): 2--12 mIU/mL (mild age-related increase)
  \item Children (pre-pubertal): <1 mIU/mL
  \item LH:FSH ratio in PCOS: >2:1 (supportive finding)
  \item Note: Always use laboratory-specific reference ranges. Interpret alongside FSH, estradiol (female), and testosterone (male).
\end{itemize}

\section*{Follow-up Tests}
\begin{itemize}
  \item FSH: Always co-measured with LH. The LH:FSH ratio and absolute levels together classify gonadal disorders.
  \item Total and free testosterone (male): Morning sample (7--10 AM). Classifies primary vs. secondary hypogonadism.
  \item Estradiol (female): Follicular phase measurement alongside LH and FSH for infertility evaluation.
  \item Progesterone (luteal phase, day 21--23 of 28-day cycle): Confirms ovulation if >3--5 ng/mL.
  \item AMH: Ovarian reserve marker, supplements FSH/LH in infertility evaluation.
  \item Prolactin: Exclude hyperprolactinemia suppressing GnRH and LH.
  \item TSH and free T4: Exclude thyroid dysfunction affecting menstrual cycles.
  \item Karyotype: For primary testicular failure or premature ovarian insufficiency.
  \item Pituitary MRI: For secondary hypogonadism (low LH/FSH with low sex steroids).
  \item Semen analysis: In male infertility with abnormal LH/FSH/testosterone.
  \item GnRH stimulation test: For precocious puberty evaluation. LH >5--8 mIU/mL after GnRH agonist (leuprolide) confirms central precocious puberty.
\end{itemize}

\section*{References}
\bibliographystyle{unsrtnat}
\bibliography{references}

\clearpage
